\section{Introduction}
\label{sec:introduction}

Machine-learning malware detectors operate under partial observability. A
portable executable (PE) can be inspected statically before execution, providing
immediate features from its headers, sections, imports, strings, and byte
statistics. This early evidence supports high-throughput screening, but packing,
obfuscation, source shift, and novel implementation choices can weaken learned
associations \cite{ucci2019survey,gibert2020rise,arp2022dos}. Dynamic analysis
reveals API calls and other behavior that a static representation may obscure,
but only after a monitored execution has produced enough events. It adds
collection cost, depends on the analysis environment, and can miss behavior that
does not manifest during observation \cite{sikorski2012practical,
sgandurra2016automated}. The operational question is therefore broader than
choosing one detector: a system must decide how to allocate heterogeneous
evidence as it becomes available.

The same allocation problem reappears during model maintenance. New labeled
software is expensive because an analyst must inspect only a small fraction of
an incoming pool. Uncertainty sampling prioritizes objects near a model boundary
\cite{cohn1996active,settles2009active}, yet it compresses the model's evidence
into a probability margin. Two samples can have the same output probability
while exhibiting different attribution structures: one may have uniformly weak
support, whereas another may combine strong positive and negative effects that
cancel. Explanation-derived structure can therefore supply a second view of
which labels may be informative. Once explanations influence sample selection,
however, they should not be accepted from a visualization alone. Their ranked
features must be tested against the fitted model, and their sensitivity to
training randomness and data partitions must be reported
\cite{warnecke2020evaluating,adebayo2018sanity}.

We organize these requirements as \emph{auditable evidence allocation under
partial observability}. The formulation has three connected stages. At the
decision stage, an early static score and sequential dynamic evidence must be
combined. At the adaptation stage, a restricted labeling budget must be assigned
to candidate samples. At the audit stage, the explanation mechanism used in
adaptation and model inspection must be subjected to controlled diagnostics. The
connection is conceptual and architectural: the same security-learning process
must allocate online evidence, offline labels, and trust in explanatory
artifacts. It does not imply that all components were evaluated on one dataset.

The first stage draws on proportional--integral--derivative (PID) control. In a
physical feedback loop, a PID controller transforms tracking error into an
actuation signal: the proportional term responds to present error, the integral
term accumulates persistent error, and the derivative term reacts to error
change \cite{astrom2001future}. We transfer this structure cautiously. Detector
disagreement replaces physical tracking error, and the controller output changes
the relative weight of static and dynamic scores. There is no physical plant or
closed-loop stability guarantee. The design is therefore \emph{control-inspired
score fusion}. This distinction is important because component ablation, rather
than the PID name, must determine whether temporal memory or a derivative term is
useful.

The latest paired evidence changes the scientific conclusion. On 1,922
MalBehavD-V1 samples with static metadata and ordered API events, a proportional
adaptive rule reduced false-positive rate (FPR) and Brier score relative to
validation-selected fixed weighting. A full PID rule obtained ranking
performance comparable to logistic stacking and a lower FPR at its selected
operating point, but it incurred a small true-positive-rate trade-off. More
importantly, PI and PID did not improve consistently over P within the examined
1--20-event horizon. We report this negative result directly and interpret it as
evidence-driven simplification: current disagreement is useful in this cohort,
whereas accumulated and derivative disagreement have not yet earned additional
complexity.

The adaptation and audit stages are evaluated on a separate, identity-audited
subset of EMBER \cite{anderson2018ember}. We combine predictive uncertainty with
two summaries of signed TreeSHAP contributions: directional entropy, which
measures the balance between positive and negative attribution mass, and
attribution variance across coordinates. Five acquisition policies are compared
under matched initial sets, pools, budgets, and outer folds. We then evaluate
TreeSHAP separately over 25 fitted models by replacing sample-specific
top-$|\mathrm{SHAP}|$ coordinates with training-derived mean or median values and
comparing the result with equal-budget random masks. Spearman and Jaccard
agreement distinguish seed sensitivity from fold sensitivity.

This component-wise design answers three research questions:

\begin{itemize}
    \item \textbf{RQ1---adaptive fusion:} On paired event-indexed static and
    dynamic evidence, how do P, PI, and PID fusion compare with average,
    max-risk, fixed weighting, and learned logistic stacking in discrimination,
    calibration, and operating-point trade-offs?
    \item \textbf{RQ2---label allocation:} Under matched training-label budgets,
    do directional entropy and SHAP variance add useful acquisition information
    beyond random and uncertainty-only selection?
    \item \textbf{RQ3---explanation audit:} Under training-derived mean and
    median replacement, do top-$|\mathrm{TreeSHAP}|$ coordinates affect fitted-
    model ranking more than equal-budget random coordinates, and how stable are
    the rankings across model seeds and data folds?
\end{itemize}

The study makes four contributions.

\begin{itemize}
    \item We formulate one auditable evidence-allocation problem that connects
    sequential decision fusion, budgeted model adaptation, and explanation
    verification while keeping their empirical estimands separate.
    \item We provide a paired comparison of simple and learned fusion rules on
    1,922 samples. Adaptive P improves FPR by 1.28 percentage points and Brier by
    0.00306 relative to fixed weighting, with paired-bootstrap intervals that
    exclude zero. The ablation also shows that I and D do not provide stable
    incremental value in the measured horizon.
    \item We evaluate explanation-guided active learning with five outer folds
    and five seeds. $U+H$ leads ROC/AP area under the learning curve (AULC),
    $U+V$ reaches the validation target with the fewest labels, and \xfull
    uses 92 fewer training labels than Random. No policy dominates every
    criterion.
    \item We test TreeSHAP rankings with controlled mean/median deletion and
    stability analysis. Top-ranked deletion produces substantially larger AUC
    decreases than random deletion, while moderate cross-model agreement exposes
    dependence on the fitted cohort.
\end{itemize}

The scope is deliberately narrower than a deployment claim. MalBehavD-V1 is
used for fusion; EMBER is used for active learning and fidelity. Their metrics
are not merged, and no active-learning--fusion interaction is claimed. Both
datasets have generic benign/malware labels, so the study does not establish
ransomware-specific or zero-day detection. The fusion comparison is development
evidence obtained after the original confirmatory analysis on the same cohort;
the active-learning and fidelity results are confirmatory under the retained
protocol. These evidence levels are visible in the tables and discussion.

The remainder of the paper is organized as follows.
Section~\ref{sec:related} reviews static/dynamic detection, score fusion, active
learning, and explanation validation. Section~\ref{sec:problem} states the
unified problem and claim boundaries. Section~\ref{sec:method} describes the
fusion, acquisition, and audit methods. Section~\ref{sec:protocol} separates the
two benchmark protocols and evidence tiers. Section~\ref{sec:results} reports
results by research question. Sections~\ref{sec:discussion} and
\ref{sec:threats} interpret the findings and threats to validity, and
Section~\ref{sec:conclusion} concludes.
